% =======================================================
% 3. ARCHITETTURA DEL SISTEMA
% =======================================================
\section{Architettura del Sistema}

\subsection{Architettura Generale}
Il sistema \textit{Second Brain} è basato su un'architettura di tipo Client-Server per garantire la separazione delle responsabilità (Separation of Concerns) tra la gestione dell'interfaccia utente (editor Markdown) e l'orchestrazione delle chiamate ai Modelli di Linguaggio di Grande Scala (LLM).
[Inserire diagramma architettura di alto livello]

\subsection{Architettura di Deployment}
[Descrizione della strategia di containerizzazione scelta (es. Docker). Specificare come vengono deployati il Frontend, il Backend e l'eventuale strato di persistenza dati/database].

\subsection{Architettura del Frontend}
Il Frontend gestisce l'interazione dell'utente con l'editor Markdown e la visualizzazione dello stato delle operazioni AI.
[Descrivere il pattern architetturale scelto per il frontend, es. MVVM, Architettura a Componenti, ecc.]

\subsubsection{Gestione dello Stato e dell'Editor}
[Descrivere come il sistema mantiene lo stato dell'editor (testo inserito, renderizzazione Markdown) e come viene garantita la sincronizzazione dei dati per l'auto-salvataggio o la persistenza locale].

\subsubsection{Gestione Interazioni AI (Widget/Pannello AI)}
[Descrivere l'architettura dei componenti responsabili di inviare il prompt, mostrare lo stato di caricamento (es. spinner o streaming del testo) e permettere l'inserimento/scarto della risposta dell'LLM].

\subsection{Architettura del Backend}
Il Backend funge da orchestratore centrale e da strato di astrazione verso i provider LLM esterni.
[Descrivere il pattern architetturale scelto per il backend, es. Architettura a Livelli (Layered), Architettura Esagonale (Ports and Adapters), ecc.]

\subsubsection{Livello di Presentazione / Controller}
[Descrizione di come il backend riceve le richieste dal client relative a operazioni testuali o di AI].

\subsubsection{Livello di Business Logic (Core LLM)}
[Descrizione del motore centrale: come il backend costruisce dinamicamente i prompt, come gestisce le logiche dei "Sei Cappelli di De Bono", del "Distant Writing" e della traduzione].

\subsubsection{Livello di Integrazione / Infrastructure}
[Descrizione di come il backend maschera le dipendenze specifiche dei provider LLM (es. OpenAI, modelli locali) e gestisce la formattazione delle risposte per il client].